\section{面向研究团队的落地清单}

\subsection{从课程到工程：如何开始搭建可复现 Agent 研发流}
\begin{figure}[H]
\centering
\includegraphics[width=0.93\textwidth]{frame-08-research-questions.jpg}
\caption{研究问题收束：能力、安全、人类参与三角}
\end{figure}
\noindent\footnotesize 画面证据：字幕约在 00:51:01--00:52:00，讲者强调系统性主题与后续研究路线。\normalsize

\begin{importantbox}{建议的最小落地路径}
\begin{enumerate}
    \item 先定义可验证任务，再讨论 Agent 架构；
    \item 先建立日志与回放，再优化策略；
    \item 先做失败分解，再追求榜单分数；
    \item 先统一评测协议，再扩大团队并行。
\end{enumerate}
\end{importantbox}

\subsection{建议的评测与治理面板}
\begin{table}[H]
\centering
\begin{tabular}{p{2.8cm}p{4.1cm}p{4.6cm}}
\toprule
\textbf{面板维度} & \textbf{核心指标} & \textbf{落地建议} \\
\midrule
任务表现 & pass rate, success@k, latency & 固定任务版本，保留历史可回放轨迹 \\
执行质量 & 工具失败率、重试率、回滚率 & 引入强制审计日志与失败分类标签 \\
安全约束 & 越权事件、敏感操作告警 & 明确权限边界与人工确认节点 \\
科学复现 & 结果方差、环境一致性、复现实验次数 & 公开配置快照与评测脚本版本 \\
\bottomrule
\end{tabular}
\caption{Agent 研究从“会跑”到“可积累”的监控面板}
\end{table}

\begin{knowledgebox}{团队协作中的关键组织设计}
研究负责人应把“可复现性”设为一等目标：  
代码提交、数据版本、模型版本、评测版本必须能形成同一条追踪链，避免“结果正确但过程不可追踪”。
\end{knowledgebox}

\begin{warningbox}{不要过早追求“全自动”}
在高风险、高不确定任务里，合理的人类介入不是倒退，而是提升系统稳健性的必要机制。  
真正目标是“可控自动化”，不是“无人值守幻想”。
\end{warningbox}

\subsection{本章小结}
落地层面最重要的不是某个框架名称，而是：任务定义、评测协议、日志治理、权限控制四件事能否同时成立。


